\section{A complete worked example: the Sierpiński topos}\label{sec:sierpinski}

Before running the generator, we compute every bridge of this text on the smallest interesting topos, so that each abstract statement has a concrete picture.

\subsection{Three presentations}

Let $\mathbf 2$ be the category with two objects $0,1$ and one non-identity arrow $u : 0\to 1$. A presheaf on $\mathbf 2$ is a pair of sets $F(1)$, $F(0)$ with a map $\res{-}{u} : F(1)\to F(0)$. Think of $F(1)$ as \q{global states}, $F(0)$ as \q{local states}, and of $\res xu$ as \q{what the local observer sees of the global state $x$}.

The \emph{Sierpiński space} $\mathbb S$ has two points $o$ (open) and $c$ (closed), and three open sets $\emptyset\subset\{o\}\subset\mathbb S$. A sheaf on $\mathbb S$ is determined by $F(\mathbb S)\to F(\{o\})$ (the value on $\emptyset$ is a singleton by gluing over the empty cover), so
\[
\Sh(\mathbb S)\ \simeq\ \PSh(\mathbf 2)\ \simeq\ \text{the category of maps of sets }X_1\to X_0.
\]
It is also the Alexandrov space of the preorder $c\le o$ on two points (\S\ref{ssec:alexandrov}), and the classifying topos of the theory of a \emph{subterminal object}, a \q{proposition}: a model in $\Set$ is a truth value. Three presentations: a small category, a space, a theory.

\subsection{Negation}

Sub-presheaves of $F = (X_1\to X_0)$ are pairs $A_1\subseteq X_1$, $A_0\subseteq X_0$ with $\res{A_1}{u}\subseteq A_0$. By Lemma~\ref{lem:neg}:
\begin{itemize}
\item $x\in X_0$ is in $\neg A$ iff $x\notin A_0$;
\item $x\in X_1$ is in $\neg A$ iff $x\notin A_1$ and $\res xu\notin A_0$;
\item $x\in X_1$ is in $\neg\neg A$ iff $\res xu\in A_0$ (the restriction lands in $A$, and then every further restriction does).
\end{itemize}
So $x\in X_1$ is in neither $\neg A$ nor $\neg\neg A$ iff ($x\in A_1$ or $\res xu\in A_0$) and $\res xu\notin A_0$, impossible since $x\in A_1$ forces $\res xu\in A_0$. Hence $\neg A\cup\neg\neg A = F$: the Sierpiński topos is \emph{De Morgan}. It is not \emph{Boolean}: for $F = \y 1 = (\{\ast\}\to\{\ast\})$ and $A = (\emptyset,\{\ast\})$, we get $\neg A = \emptyset$, so $A\cup\neg A\neq F$.

\textbf{Check with the arches.} Bridge~\ref{br:bool}: $\mathbf 2$ is not a groupoid ($u$ is not invertible). Bridge~\ref{br:dm}: $\mathbf 2$ has the initial object $0$, so it satisfies Ore. The topological arch: the closure of the open set $\{o\}$ is $\mathbb S$, which is open, and the closures of $\emptyset$ and $\mathbb S$ are themselves: $\mathbb S$ is extremally disconnected. Three readings, one answer.

\subsection{Points}

$\Kar(\mathbf 2) = \mathbf 2$, since the only idempotents are identities. By Bridge~\ref{br:pts}, the topos has exactly two essential points, evaluation at $1$ ($F\mapsto X_1$) and at $0$ ($F\mapsto X_0$), and exactly one natural transformation between them, from evaluation at $1$ to evaluation at $0$, given by $\res{-}{u}$. Topologically, these are the stalks at $c$ and at $o$; logically, the two models of the theory of a proposition, \q{false} and \q{true}, with the unique homomorphism false $\to$ true.

\subsection{Subtopoi and topologies}

Let us find all Grothendieck topologies on $\mathbf 2$. The sieves on $0$ are $M_0 = \{\id_0\}$ and $\emptyset$; the sieves on $1$ are $M_1 = \{\id_1,u\}$, $S = \{u\}$ and $\emptyset$. Pulling back along $u$ sends $S$ to $M_0$ and $\emptyset$ to $\emptyset$. A topology always contains $M_0$ and $M_1$. One checks from the axioms that there are exactly four topologies:
\begin{enumerate}
\item the \emph{trivial} one (only $M_0$, $M_1$ cover): sheaves are all presheaves, the whole topos;
\item $S$ covers $1$: a sheaf has $X_1\to X_0$ bijective, so sheaves $\simeq\Set$ via $X_0$ --- the open point $\{o\}$;
\item $\emptyset$ covers $0$: a sheaf has $X_0$ a singleton, so sheaves $\simeq\Set$ via $X_1$ --- the closed point $\{c\}$;
\item everything covers: only the terminal presheaf is a sheaf --- the empty subspace.
\end{enumerate}
Indeed, if $\emptyset$ covers $0$ and $S$ covers $1$, then $\emptyset$ covers $1$ by transitivity (its pullback along the only arrow of $S$ is $\emptyset$, which covers $0$), and then everything covers; and if $\emptyset$ covers $1$, its pullback $\emptyset$ covers $0$. So Bridge~\ref{br:sub} says that $\Sh(\mathbb S)$ has exactly four subtopoi, which are the four subspaces $\mathbb S$, $\{o\}$, $\{c\}$, $\emptyset$. The order is reversed: the trivial topology (fewest covers) gives the largest subtopos.

\subsection{What the example teaches}

Every bridge of this text was visible here: a logical property (De Morgan) read on a category (Ore), on a space (extremally disconnected) and by hand (Lemma~\ref{lem:neg}); points read as idempotents, as stalks and as models; subtopoi read as topologies and as subspaces. In the following sections the shapes are less small and the readings less obvious, and that is where the bridges start to be useful.

\begin{exercise}\label{exo:three}
Do the same for the chain $\mathbf 3 = \{0<1<2\}$: show that $\PSh(\mathbf 3)$ is De Morgan and not Boolean, has three essential points, and find its subtopoi. (There are eight; compare with the subspaces of the Alexandrov space of $\mathbf 3^{\op}$.)
\end{exercise}
